% ===================================================== failure recovery ====
\section{Failure recovery strategy}

\subsection{The degradation ladder}

The governing principle: \emph{an emergency system that stops dispatching
because a microservice is unhealthy has failed worse than one that dispatches
on a slightly stale estimate.} Each rung below is implemented and exercised.

\begin{center}\footnotesize
\begin{tabular}{@{}rlp{56mm}l@{}}
\toprule
& \textbf{Rung} & \textbf{Trigger and behaviour} & \textbf{Quality} \\
\midrule
1 & Exact solver + live routing & Normal operation & Optimal \\
2 & Exact solver + cached travel times & Routing engine slow; cache serves & Near-optimal \\
3 & Greedy + great-circle estimates & Circuit breaker open; conservative fallback distances & Good \\
4 & Region-local greedy, last-known state & Event log partition unavailable; worker uses local snapshot & Degraded but safe \\
5 & Human dispatcher console & Decision plane unavailable; operators dispatch manually against live resource state & Manual \\
\bottomrule
\end{tabular}
\end{center}

\begin{keybox}
\textbf{The fallback is pessimistic by construction.} The great-circle
estimate applies a $1.15\times$ penalty and a speed derate, so degraded
estimates are \emph{longer} than reality rather than shorter. Over-promising
an arrival time in an emergency is the more dangerous error: it causes the
dispatcher to stop looking for alternatives.
\end{keybox}

\subsection{Failure modes and responses}

\begin{center}\footnotesize
\begin{tabular}{@{}p{34mm}p{44mm}p{50mm}@{}}
\toprule
\textbf{Failure} & \textbf{Blast radius} & \textbf{Recovery} \\
\midrule
Optimiser worker crash & One region shard & Leases expire and are reaped; partition reassigned; state rebuilt from compacted topic + database. RTO $<$ 60\,s \\
Routing engine outage & Estimate quality & Circuit breaker opens after 5 failures; fallback serves; half-open probe restores automatically \\
Database primary loss & Command path & Synchronous replica promotion; RPO $\approx$ 0 because commits require in-sync acknowledgement \\
Event log broker loss & Throughput & RF=3 with \texttt{min.insync.replicas}=2; producers retry idempotently \\
Poison message & One partition & Bounded retry, then dead-letter with full context; partition continues \\
Network partition between regions & Cross-region borrowing only & Shards remain independently operational (static stability); borrowing degrades to none \\
Cache cluster loss & Latency & Cache-aside means correctness is unaffected; load falls back to the routing engine, protected by the breaker \\
\bottomrule
\end{tabular}
\end{center}

\subsection{Consistency and delivery semantics}

Transport is at-least-once; \emph{effects} are exactly-once. That combination
is achieved with two standard patterns whose absence causes subtle corruption:

\begin{itemize}
\item \textbf{Transactional outbox.} The assignment row and the event
      announcing it are written in one transaction; a relay publishes
      afterwards. This removes the dual-write race in which a resource is
      marked busy but the assignment event never reaches the crew, or the
      reverse.
\item \textbf{Idempotency keys.} Every command carries one; an effects table
      deduplicates. Redelivery after a broker retry is therefore a no-op
      rather than a second ambulance.
\end{itemize}

\subsection{Validating resilience}

Chaos injection is part of the evaluation harness, not a separate exercise: the
reference run includes road closures, vehicle failures, hospital saturation and
a mid-surge routing-engine outage. Section~\ref{sec:eval} reports how many
travel-time estimates were served in degraded mode and how many times the
circuit breaker tripped. In production this extends to scheduled game days,
periodic restore drills, and shadow-mode replay of historical event streams
against candidate policy versions.

% ======================================================== scalability ======
\section{Scalability strategy}

\subsection{Geographic sharding as the scaling unit}

\begin{center}
\begin{tikzpicture}[node distance=4mm]
\node[bus, minimum width=112mm] (log) {Event log: 256 partitions, key = \texttt{region\_id}};
\foreach \i/\n in {0/A, 1/B, 2/C, 3/D} {
  \node[svc, below=9mm of log, xshift={(\i-1.5)*29mm}, minimum width=26mm] (w\i)
    {Solver worker \n\\\tiny owns shard, exclusive writer};
  \draw[fl] ([xshift={(\i-1.5)*29mm}]log.south) -- (w\i.north);
}
\node[store, below=8mm of w1, xshift=14mm, minimum width=52mm] (db)
  {Shared state: PostgreSQL (CAS) + Redis (leases)\\\tiny consulted only for cross-shard borrowing};
\foreach \i in {0,1,2,3} { \draw[fl,dashed] (w\i) -- (db); }
\end{tikzpicture}
\end{center}

One region shard maps to one event-log partition and one solver worker holding
exclusive write ownership. Shards run in parallel with no shared state and no
distributed locking on the hot path. The national fleet scales by adding
shards, not by growing one matrix quadratically.

\begin{center}\footnotesize
\begin{tabular}{@{}lp{46mm}p{52mm}@{}}
\toprule
\textbf{Dimension} & \textbf{Bottleneck} & \textbf{Scaling response} \\
\midrule
Incident rate & Ingestion CPU & Stateless horizontal scale behind the gateway; over-partition the log (256) so future growth needs no repartitioning \\
Fleet size & Solver matrix $O(n^3)$ & Per-shard candidate and slot caps; the spatial index bounds candidates by geometry, not by fleet size \\
Concentrated disaster & One hot shard & Sub-shard split: the region is divided at cell level and additional workers claim the sub-keys \\
Read/query load & Database replicas & CQRS: command side on PostgreSQL, read models materialised in Redis and ClickHouse \\
Cross-region coordination & Central coordinator & Hierarchical: local solvers plus a coarse national transfer problem solved at a longer cadence \\
\bottomrule
\end{tabular}
\end{center}

\begin{decbox}
\dec{Autoscale on consumer lag, not CPU.}{Lag is the metric that actually
correlates with dispatch delay. A solver worker can be at 40\% CPU and
minutes behind because it is blocked on I/O, and CPU-based autoscaling would
never notice. KEDA-style lag-driven scaling reacts to the thing the citizen
experiences.}
\end{decbox}

\subsection{Hot-partition handling}

A cyclone concentrates demand into two or three regions. Uniform partitioning
would leave one worker saturated while twenty idle. Detection is by lag and
slot-count per shard; the response is to split the region's key space at H3
cell granularity and let additional workers claim sub-keys, with the
reservation lease continuing to guarantee conflict-freedom across the split.

% ======================================================== performance ======
\section{Performance considerations}

\subsection{Latency budgets and measured results}

\begin{center}\footnotesize
\begin{tabular}{@{}lrrl@{}}
\toprule
\textbf{Path} & \textbf{Budget} & \textbf{Measured (p99)} & \textbf{Notes} \\
\midrule
Tier-1 admission decision & 200\,ms & \textbf{2.4\,ms} & Reference implementation, pure Python, single core \\
Tier-2 shard epoch & 400\,ms & \textbf{36\,ms} & Exact solve + LNS across all shards \\
Travel-time lookup (cached) & 1\,ms & --- & 97--98\% hit rate observed \\
End-to-end accept $\rightarrow$ commit & 200\,ms & --- & Sum of stages 1--7 \\
\bottomrule
\end{tabular}
\end{center}

\noindent The headroom is deliberate. A pure-Python reference implementation
running two orders of magnitude inside budget means a production
implementation in Go or Rust has room for fleets far larger than the one
simulated, and room for the budget to be consumed by network hops that the
in-process reference does not pay.

\subsection{Where the performance actually comes from}

\begin{itemize}
\item \textbf{Bounded candidate sets.} The spatial index caps candidates
      geometrically, so Tier-1 cost is independent of national fleet size.
      This is the single most important performance property: without it,
      admission latency degrades as the system succeeds and grows.
\item \textbf{Contraction hierarchies} in the routing engine give
      microsecond-scale path queries; AEGIS additionally caches at cell
      granularity so most queries never reach it.
\item \textbf{Cached free-flow, overlaid live.} The cache stores free-flow
      travel time \textemdash{} a property of the road graph \textemdash{} and
      applies closure and weather multipliers at read time. This is why a
      coarse key and a ten-minute TTL are safe; baking the overlay in would
      make both unsafe.
\item \textbf{Binary serialisation} (Protobuf with a schema registry) and
      batched database writes behind a connection pooler.
\item \textbf{Load shedding by priority.} Under saturation, telemetry is shed
      before incident intake. The last thing to degrade is the ability to
      accept an emergency call.
\end{itemize}

% ========================================================== database =======
\section{Database design}

\subsection{Polyglot persistence, justified per store}

\begin{center}\footnotesize
\begin{tabular}{@{}lp{40mm}p{56mm}@{}}
\toprule
\textbf{Store} & \textbf{Holds} & \textbf{Why this store} \\
\midrule
PostgreSQL + PostGIS & Incidents, resources, assignments, outbox, audit & Multi-entity transactional integrity is non-negotiable for reservations; PostGIS gives native geography indexing. A pure key-value store cannot express the all-or-nothing group reservation \\
Redis & Live positions (GEO sets), leases, ETA cache, dedup keys, rate limits & Sub-millisecond access to state that changes on every telemetry ping; TTL semantics are exactly what leases need \\
ClickHouse / TimescaleDB & Telemetry, metrics, decision journal & Telemetry write amplification would destroy the command database; columnar storage suits the analytical query shape \\
Object store (WORM) & Sealed audit records & Immutable retention for post-incident inquiry; cheap at long horizons \\
\bottomrule
\end{tabular}
\end{center}

\subsection{Core schema}

\begin{center}\footnotesize
\begin{tabular}{@{}lp{86mm}@{}}
\toprule
\textbf{Table} & \textbf{Key columns, keys and indexes} \\
\midrule
\texttt{incident} &
\texttt{incident\_id} PK, \texttt{region\_id}, \texttt{geog GEOGRAPHY(Point)},
\texttt{severity}, \texttt{affected}, \texttt{window\_s}, \texttt{hazard},
\texttt{state}, \texttt{version}, \texttt{reported\_at}.
GIST on \texttt{geog}; BRIN on \texttt{reported\_at}; partial B-tree on
\texttt{state} where open. Partitioned monthly by \texttt{reported\_at},
sub-partitioned by \texttt{region\_id} \\
\addlinespace
\texttt{resource} &
\texttt{resource\_id} PK, \texttt{type}, \texttt{region\_id}, \texttt{geog},
\texttt{capability}, \texttt{capacity}, \texttt{occupied}, \texttt{state},
\texttt{version}. GIST on \texttt{geog}; partial index on
\texttt{state='AVAILABLE'} \\
\addlinespace
\texttt{assignment} &
\texttt{assignment\_id} PK, \texttt{incident\_id} FK, \texttt{resource\_id} FK,
\texttt{state}, \texttt{eta\_s}, \texttt{cost}, \texttt{generation},
\texttt{supersedes}, \texttt{rationale JSONB}.
\textbf{Partial unique index on \texttt{resource\_id} where state is live}
\textemdash{} the database itself refuses to record a double-booking \\
\addlinespace
\texttt{assignment\_event} &
Append-only event stream per assignment aggregate. The state table is a
projection; this is the source of truth and the replay substrate \\
\addlinespace
\texttt{outbox} &
\texttt{topic}, \texttt{key}, \texttt{payload}, \texttt{published\_at}.
Written in the same transaction as the state change; drained by a relay \\
\bottomrule
\end{tabular}
\end{center}

\begin{keybox}
\textbf{The partial unique index is the point.} Application-level conflict
prevention is necessary but not sufficient; a database constraint that makes
the invalid state \emph{unrepresentable} converts a possible silent corruption
into a loud, immediate constraint violation. Defence in depth, at negligible cost.
\end{keybox}

\subsection{Event sourcing for assignments only}

Assignments are event-sourced; incidents and resources are not. The asymmetry
is deliberate: assignments are the entity whose history must be reconstructible
in an inquiry ("why did unit 47 go to the wrong address?"), and whose state
transitions carry legal weight. Event-sourcing everything would impose
projection complexity where a mutable row with a version column is sufficient.

% ========================================================== caching ========
\section{Caching strategy}

\begin{center}\footnotesize
\begin{tabular}{@{}lp{34mm}p{30mm}p{34mm}@{}}
\toprule
\textbf{Layer} & \textbf{Contents} & \textbf{Policy} & \textbf{Invalidation} \\
\midrule
L1 in-process & Region polygons, hospital metadata, policy weights & TTL + version & Broadcast on \texttt{config.updates} \\
L2 Redis & Travel-time matrix, resource snapshots, incident read models & Cache-aside, 10\,min TTL & Partition-scoped version bump \\
Write-through & Resource live state & Immediate & Written with the state change \\
Negative cache & "no route exists" results & Short TTL & Prevents repeated futile routing calls \\
\bottomrule
\end{tabular}
\end{center}

\subsection{Surgical invalidation}

When a bridge closes, only the travel times touching the affected cells must be
invalidated. Each cache entry records the \emph{partition version} it was
written under; a version bump makes every affected entry a miss with no scan
and no delete pass, while unrelated routes stay hot. A nationwide flush during
a cyclone would collapse into a routing-engine stampede exactly when the system
is busiest. A unit test asserts that an unrelated route remains cached across a
closure.

\subsection{Stampede protection}

Two mechanisms, both necessary:

\begin{itemize}
\item \textbf{Probabilistic early expiry (XFetch).} Under surge, thousands of
      concurrent requests miss the same key at the same instant. Staggering
      refreshes before hard expiry spreads the recomputation.
\item \textbf{Single-flight.} An in-flight set collapses duplicate
      recomputations of the same key to one, so only one caller pays the
      routing call.
\end{itemize}

% ================================================== event processing =======
\section{Event processing strategy}

\begin{center}\footnotesize
\begin{tabular}{@{}lp{88mm}@{}}
\toprule
\textbf{Property} & \textbf{Design} \\
\midrule
Ordering & Total order per partition (region). Cross-region ordering is not required and is not paid for \\
Delivery & At-least-once transport; exactly-once effects via idempotency keys \\
Backpressure & Bounded queues; producers block rather than buffer without limit; load shed by priority class \\
Retry & Bounded exponential retry in-partition, then dead-letter with full context and replay tooling \\
Fairness & Partitions polled round-robin with a bounded batch, so one saturated region cannot starve the rest of the country \\
Schema & Protobuf with a registry; backward-compatible evolution enforced in CI \\
Replay & The log plus the decision journal allow any historical period to be re-run against a candidate policy version in shadow mode \\
\bottomrule
\end{tabular}
\end{center}

\begin{keybox}
\textbf{Replay is a first-class capability, not a debugging afterthought.}
It is how a proposed change to the triage weights is evaluated before it
touches a real dispatch: re-run last month's actual event stream under the new
policy, compare outcome distributions, and only then promote. A tuning change
in this domain is a clinical change, and deserves the same evidence standard.
\end{keybox}

% ========================================================== security =======
\section{Security considerations}

\begin{center}\footnotesize
\begin{tabular}{@{}lp{88mm}@{}}
\toprule
\textbf{Concern} & \textbf{Control} \\
\midrule
Service identity & mTLS between all services via a service mesh; workload identity, no shared secrets \\
Human access & OIDC with short-lived tokens; RBAC combined with ABAC scoped to region, so a district operator cannot retask another division's fleet \\
Casualty data & Field-level encryption with envelope keys in a KMS; logs carry pseudonymous incident ids only \\
Audit & Append-only decision journal, signed and sealed to WORM storage; immutable for the statutory retention period \\
\textbf{Malicious incident injection} & The highest-value attack: fabricated emergencies that drain a district's fleet before a real attack. Mitigated by reporter reputation scoring, cross-verification of clustered reports, per-source rate limits, and anomaly detection on arrival-rate deviation from forecast. Suspicious reports are dispatched at reduced weight rather than dropped \\
Resource-drain abuse & Per-caller quotas; coverage floors that the optimiser will not breach regardless of incoming priority \\
Secrets & Vault-managed with automatic rotation; least privilege per service \\
Network & Segmented; the decision plane is not internet-reachable \\
Compliance & Health data handled to HIPAA/GDPR-equivalent standards; data residency enforced per region \\
\bottomrule
\end{tabular}
\end{center}

\begin{decbox}
\dec{Suspicious reports are down-weighted, not rejected.}{A false negative
here is a real emergency ignored because the reporter had no history. The
scoring model can express reduced confidence as a lower vulnerability and
scale prior; it should not express it as a dropped call.}
\end{decbox}

% ==================================================== observability =======
\section{Monitoring and observability}

Infrastructure metrics are necessary but not sufficient. The signals that
matter here are clinical, not computational.

\begin{center}\footnotesize
\begin{tabular}{@{}lp{78mm}@{}}
\toprule
\textbf{Class} & \textbf{Signals} \\
\midrule
Domain SLIs & Time-to-first-assignment; first-arrival latency by severity; deadline-miss rate; unassigned incident count and age; reassignment churn rate; coverage gap per region \\
Solver health & Objective delta per epoch; solve time; matrix dimensions; budget exhaustion count; LNS improvement rate \\
Resource health & Utilisation by type and region; lease reclaim rate; commit-denied rate; vehicles out of service \\
Dependency health & Circuit-breaker state; degraded-estimate rate; cache hit rate; consumer lag per partition \\
Standard & RED for services, USE for nodes, error budgets and burn-rate alerts \\
\bottomrule
\end{tabular}
\end{center}

\subsection{Tracing and the decision journal}

\texttt{incident\_id} propagates as OpenTelemetry baggage, so a single trace
spans call intake to hospital handover. Alongside it, every commit writes a
\textbf{decision record}:

\begin{center}\footnotesize
\begin{tabular}{@{}ll@{}}
\toprule
\textbf{Field} & \textbf{Contents} \\
\midrule
\texttt{candidates\_considered} & Size of the evaluated candidate set \\
\texttt{rejected\_reasons} & Histogram: type mismatch, capability below minimum, beyond reach radius, unavailable \\
\texttt{cost\_breakdown} & Every term: travel, lateness, over-qualification, scarcity, coverage, continuity, switching \\
\texttt{priority\_breakdown} & Per-term triage contributions and policy version \\
\texttt{runner\_up\_cost} & What the next-best option would have cost \\
\texttt{objective\_before/after} & The improvement this epoch produced \\
\bottomrule
\end{tabular}
\end{center}

\begin{keybox}
This record is the substrate for three otherwise-impossible things: the
"why this unit?" panel in the dispatcher console, the post-incident inquiry,
and offline policy evaluation. It is written on the commit path, not
reconstructed later, because a reconstruction is an argument and a record is
evidence.
\end{keybox}

\subsection{Human override}

Dispatchers can override any assignment. Overrides are recorded with the
system's own recommendation and the operator's stated reason, and the
divergence rate between human and system choices is itself a monitored signal:
a rising divergence rate in one region is evidence that the model's weights no
longer match conditions on the ground.
